\documentclass[11pt,a4paper]{article}
\usepackage[utf8]{inputenc}
\usepackage[T1]{fontenc}
\usepackage{geometry}
\usepackage{booktabs}
\usepackage{longtable}
\usepackage{float}
\usepackage{array}
\usepackage{amsmath}
\usepackage{graphicx}
\usepackage{hyperref}
\usepackage{xcolor}
\geometry{margin=1in}

\title{Custom-WikiCS Graph Analysis -- Part 8\\[4pt]
\large Metric Expansion on the Revised Undirected Structural Branch}
\author{Sahin Uregil}
\date{\today}

\begin{document}
\maketitle

\section{Introduction}
This report extends the revised undirected structural branch introduced in
Report 7. The underlying experimental branch does not change: the same
undirected train/test split, the same ten-model family, and the same
retrained Node2Vec, GCN, and GraphSAGE artifacts are reused. What changes here
is the \emph{measurement layer}.

Two follow-up evaluation passes were added:

\begin{itemize}
    \item \textbf{Eval 3.1}: ranking metrics at a smaller cutoff together with
          link-discrimination metrics
    \item \textbf{Eval 3.2}: relative-improvement analysis across the key model
          families
\end{itemize}

The goal of this report is therefore not to replace Report 7, but to refine
its conclusions. In particular, it asks three questions:

\begin{itemize}
    \item What changes when ranking quality is measured at \textbf{@10} rather
          than through larger-cutoff metrics?
    \item How do the models compare under \textbf{ROC-AUC} and \textbf{PR-AUC}
          when viewed as link-discrimination systems?
    \item How large are the improvements from hybridization, pooling, and
          power-law correction when expressed as \textbf{percentage gains}?
\end{itemize}

\section{Evaluation Additions}

\subsection{Smaller-Cutoff Ranking Metrics}
Report 7 already used Hit@10, but several other ranking metrics were still
reported only at the larger effective cutoff inherited from the original
evaluation family. Eval 3.1 therefore adds:

\begin{itemize}
    \item NDCG@10
    \item Recall@10
    \item MAP@10
\end{itemize}

The project still keeps plain MRR unchanged rather than introducing a truncated
``MRR@10'', because standard MRR is already well-defined and easier to compare
against the earlier report line.

One structural note is important here: the current benchmark still uses one
positive target per query. Because of that, \textbf{Recall@10 is numerically
equivalent to Hit@10}. It is retained for completeness because it belongs to
the requested metric family, but its interpretation is the same in this
single-positive setting.

\subsection{ROC-AUC and PR-AUC}
Eval 3.1 also adds two link-discrimination metrics for the full model family:

\begin{itemize}
    \item ROC-AUC
    \item PR-AUC
\end{itemize}

The procedure is:

\begin{itemize}
    \item use the revised-branch undirected test edges as positives
    \item sample the same number of undirected non-edges as negatives
    \item score each edge using the same similarity layer used by the ranking
          evaluation
\end{itemize}

For the reciprocal-rank-fusion model, the underlying score is directional, so
an undirected edge score is defined as the mean of the $u \rightarrow v$ and
$v \rightarrow u$ fusion scores.

\section{Cutoff-10 and AUC Results}

\begin{table}[H]
\centering
\caption{Eval 3.1 extension metrics on the revised branch. Ranking metrics use 86,446 expanded evaluation pairs. ROC-AUC and PR-AUC use 43,223 positive and 43,223 sampled negative undirected edges.}
\setlength{\tabcolsep}{4pt}
\renewcommand{\arraystretch}{1.12}
\resizebox{\textwidth}{!}{
\begin{tabular}{@{}l*{6}{c}@{}}
\toprule
\textbf{Model} &
\textbf{H@10} &
\textbf{NDCG@10} &
\textbf{Recall@10} &
\textbf{MAP@10} &
\textbf{ROC-AUC} &
\textbf{PR-AUC} \\
\midrule
BGE pooled + Node2Vec (corrected) & 8.67\% & 0.0409 & 8.67\% & 0.0273 & 0.9563 & 0.9559 \\
BGE-M3 + Node2Vec (corrected) & 8.67\% & 0.0413 & 8.67\% & 0.0278 & 0.9582 & 0.9579 \\
BGE pooled + Node2Vec (uncorrected) & 7.83\% & 0.0356 & 7.83\% & 0.0229 & 0.9710 & 0.9716 \\
BGE-M3 + Node2Vec (uncorrected) & 7.83\% & 0.0356 & 7.83\% & 0.0229 & 0.9710 & 0.9716 \\
Node2Vec (uncorrected) & 7.81\% & 0.0355 & 7.81\% & 0.0229 & 0.9717 & 0.9714 \\
Node2Vec (corrected) & 7.81\% & 0.0355 & 7.81\% & 0.0228 & 0.9717 & 0.9714 \\
BGE-M3 Only & 7.64\% & 0.0406 & 7.64\% & 0.0298 & 0.8851 & 0.8886 \\
GCN + BGE-M3 (rank fusion) & 7.63\% & 0.0355 & 7.63\% & 0.0234 & 0.9356 & 0.9363 \\
GCN Only & 3.46\% & 0.0161 & 3.46\% & 0.0107 & 0.8954 & 0.8946 \\
GraphSAGE Only & 1.22\% & 0.0060 & 1.22\% & 0.0041 & 0.8691 & 0.8561 \\
\bottomrule
\end{tabular}
}
\end{table}

\section{What the New Metrics Change}

\subsection{The Top-of-List Story Still Favors Corrected Hybrids}
The smaller-cutoff results preserve the main ranking conclusion from Report 7:
the strongest models for recommendation-oriented top-list behavior are still
the \textbf{corrected hybrid variants}.

The top line is:

\begin{itemize}
    \item \textbf{BGE pooled + Node2Vec (corrected)} remains best on Hit@10
          with $8.6736\%$
    \item \textbf{BGE-M3 + Node2Vec (corrected)} is essentially tied on Hit@10
          and is slightly stronger on NDCG@10 and MAP@10
\end{itemize}

So the previous interpretation survives:

\begin{itemize}
    \item pooled corrected hybrid is the best top-line Hit@10 model
    \item non-pooled corrected hybrid remains the strongest corrected model on
          several secondary ranking metrics
\end{itemize}

\subsection{AUC and PR-AUC Tell a Different Story}
The AUC layer produces a more nuanced result. Under ROC-AUC and PR-AUC, the
best models are no longer the corrected hybrids. Instead, the strongest
link-discrimination scores come from:

\begin{itemize}
    \item \textbf{Node2Vec (uncorrected)}: ROC-AUC $= 0.971666$, PR-AUC $= 0.971390$
    \item \textbf{Node2Vec (corrected)}: ROC-AUC $= 0.971666$, PR-AUC $= 0.971376$
    \item the uncorrected hybrid variants, which sit very close behind
\end{itemize}

This creates an important distinction:

\begin{itemize}
    \item \textbf{corrected hybrids are best for recommendation-oriented
          top-ranked output}
    \item \textbf{pure structural or lightly fused uncorrected models are best
          for global edge discrimination}
\end{itemize}

This is not a contradiction. AUC and PR-AUC measure broad separation quality
across many scored pairs, while Hit@10, NDCG@10, and MAP@10 care much more
about what appears near the top of the list.

\subsection{BGE-M3 Alone Has an Interesting Mixed Profile}
BGE-M3 Only remains relatively weak on raw edge-recovery breadth, but it is not
poor at every smaller-cutoff metric. In particular:

\begin{itemize}
    \item its Hit@10 is below the top corrected hybrids
    \item but its NDCG@10 and MAP@10 are noticeably stronger than its broad
          ranking profile would suggest
\end{itemize}

This indicates that semantic-only retrieval still puts a meaningful amount of
useful signal near the top of the list, even though it remains more popularity
driven than the stronger hybrids.

\section{Relative Improvement Analysis}

\subsection{Method}
Eval 3.2 converts the metric tables into percentage-gain comparisons. For
higher-is-better metrics, improvement is defined as:

\[
100 \cdot \frac{\text{candidate} - \text{baseline}}{\text{baseline}}
\]

For lower-is-better metrics such as Popularity Lift and Gini Index, the sign is
reversed so that lower bias still appears as a positive improvement:

\[
100 \cdot \frac{\text{baseline} - \text{candidate}}{\text{baseline}}
\]

\subsection{Focused Comparison Table}

\begin{table}[H]
\centering
\caption{Relative improvements (\%) from Eval 3.2. Positive values indicate improvement under the appropriate metric direction.}
\setlength{\tabcolsep}{4pt}
\renewcommand{\arraystretch}{1.12}
\resizebox{\textwidth}{!}{
\begin{tabular}{@{}l*{11}{c}@{}}
\toprule
\textbf{Comparison} &
\textbf{H@10} &
\textbf{NDCG@10} &
\textbf{Recall@10} &
\textbf{MAP@10} &
\textbf{ROC-AUC} &
\textbf{PR-AUC} &
\textbf{Pop. Lift} &
\textbf{Gini} &
\textbf{Agg. Div.} &
\textbf{Novelty} &
\textbf{ILD} \\
\midrule
Flagship hybrid vs text-only & 14.60 & 0.77 & 13.55 & -8.46 & 8.04 & 7.58 & 31.48 & 20.53 & 3.86 & 3.09 & 17.75 \\
Flagship hybrid vs graph-only & 11.03 & 15.22 & 11.02 & 19.26 & -1.58 & -1.60 & 0.54 & 8.29 & 3.28 & 1.62 & -4.30 \\
Pooling effect (corrected) & 0.05 & -0.95 & 0.05 & -1.81 & -0.19 & -0.22 & -1.68 & 0.41 & 0.51 & -0.03 & 1.31 \\
Pooling effect (uncorrected) & 0.01 & -0.04 & 0.01 & -0.09 & -0.00 & -0.00 & 0.02 & -0.17 & 0.20 & 0.01 & 0.03 \\
Correction effect (Node2Vec) & -0.04 & -0.12 & -0.04 & -0.21 & -0.00 & -0.00 & 0.00 & 0.00 & 0.00 & 0.00 & 0.00 \\
Correction effect (hybrid) & 10.66 & 16.00 & 10.66 & 21.12 & -1.32 & -1.40 & 2.38 & 10.38 & 0.93 & 1.74 & -5.30 \\
Correction effect (pooled hybrid) & 10.70 & 14.95 & 10.70 & 19.04 & -1.51 & -1.61 & 0.71 & 10.90 & 1.25 & 1.70 & -4.08 \\
\bottomrule
\end{tabular}
}
\end{table}

\section{Interpretation of the Relative Gains}

\subsection{Hybridization Clearly Beats Text-Only}
Using the pooled corrected hybrid as the flagship model, the gain over pure
BGE-M3 is clear:

\begin{itemize}
    \item Hit@10 improves by about \textbf{14.6\%}
    \item Recall@10 improves by about \textbf{13.6\%}
    \item ROC-AUC improves by about \textbf{8.0\%}
    \item PR-AUC improves by about \textbf{7.6\%}
    \item Popularity Lift improves by about \textbf{31.5\%}
    \item Gini improves by about \textbf{20.5\%}
    \item Intra-List Distance improves by about \textbf{17.8\%}
\end{itemize}

There is one counterpoint: MAP@10 becomes worse by about $8.5\%$. So the best
hybrid is not uniformly better in every precision-sensitive small-cutoff
measure. Still, the total balance of gains strongly favors the hybrid.

\subsection{Hybridization Beats Graph-Only at the Top of the List, but Not on AUC}
Against Node2Vec (uncorrected), the flagship hybrid gains:

\begin{itemize}
    \item about \textbf{11.0\%} on Hit@10
    \item about \textbf{15.2\%} on NDCG@10
    \item about \textbf{19.3\%} on MAP@10
\end{itemize}

But it loses:

\begin{itemize}
    \item about \textbf{1.58\%} on ROC-AUC
    \item about \textbf{1.60\%} on PR-AUC
\end{itemize}

This is one of the most important findings in the report. It shows that the
strongest recommendation-style model is not the same as the strongest
global-discrimination model.

\subsection{Pooling Helps Only a Little}
The pooling comparison is intentionally small, and the results confirm that.
For the corrected branch, pooling produces:

\begin{itemize}
    \item a tiny positive Hit@10 effect
    \item a tiny positive Recall@10 effect
    \item a small Aggregate Diversity gain
    \item but slightly worse NDCG@10, MAP@10, ROC-AUC, and PR-AUC
\end{itemize}

The uncorrected pooling comparison is even smaller and is effectively
negligible. So the correct interpretation is:

\begin{itemize}
    \item pooling is not useless
    \item but its effect is subtle rather than transformative
\end{itemize}

\subsection{Correction Matters for Hybrids, Not for Pure Node2Vec}
The Node2Vec-only correction comparison is almost perfectly flat. In practical
terms, the revised branch suggests that correction changes almost nothing for
pure Node2Vec.

The hybrid families are different. For both the non-pooled and pooled hybrid
branches, correction yields:

\begin{itemize}
    \item about \textbf{10--11\%} gain on Hit@10 / Recall@10
    \item about \textbf{15--16\%} gain on NDCG@10
    \item about \textbf{19--21\%} gain on MAP@10
\end{itemize}

At the same time, correction hurts:

\begin{itemize}
    \item ROC-AUC
    \item PR-AUC
    \item Intra-List Distance
\end{itemize}

So again the pattern is consistent: correction sharpens top-of-list
recommendation behavior for hybrids, but weakens broad edge-discrimination
behavior.

\section{Updated Conclusions}
After the metric expansion, the revised branch story becomes sharper rather
than weaker.

\begin{itemize}
    \item \textbf{The corrected hybrid family is still the best direction for
          recommendation-oriented top-ranked output.}
    \item \textbf{Pure structural models remain strongest under ROC-AUC and
          PR-AUC.}
    \item \textbf{Pooling helps a little, but its effect is subtle.}
    \item \textbf{Power-law correction matters much more for hybrids than for
          pure Node2Vec.}
\end{itemize}

The most defensible final interpretation is therefore:

\begin{quote}
The revised branch supports a two-lens conclusion. If the objective is global
link discrimination, pure structural methods remain strongest. If the objective
is recommendation quality near the top of the ranked list, corrected hybrid
semantic-structural models are the best overall choice.
\end{quote}

That distinction is important for the project because the actual product goal is
not abstract link recovery in isolation. It is article recommendation and,
eventually, learning-path generation. Under that objective, the corrected
hybrid family remains the most convincing final system direction.

\end{document}
